\documentclass[12pt]{article}
\usepackage[T1]{fontenc}
\usepackage[utf8]{inputenc}
\usepackage{lmodern}
\usepackage{textcomp}
\usepackage{enumitem}
\usepackage{geometry}
\geometry{margin=1in}
\begin{document}
\begin{center}
Answer Key - Religious Studies - K-12 Level Assessment 21 \
\small (Answers are ordered exactly as in the question paper.)
\end{center}

\section*{Multiple Choice}
\begin{enumerate}[label=\arabic*.]
\item \textbf{Answer:} C
\\ \textit{Options:} A) Guru Nanak; B) Guru Arjan; C) Guru Hargobind; D) Guru Ram Das
\\ \textit{Note:} Guru Hargobind, the sixth Sikh Guru, is known for donning two swords, representing Miri and Piri, which symbolize the dual authority of spiritual leadership and temporal power in Sikhism.
\item \textbf{Answer:} A
\\ \textit{Options:} A) 48; B) 60; C) 12; D) 36
\\ \textit{Note:} Jaina devotees practice meditation or reflection for 48 minutes each day as a means to cultivate mindfulness and spiritual growth in accordance with their religious beliefs and traditions.
\item \textbf{Answer:} C
\\ \textit{Options:} A) The Battle of Uhud; B) The Battle of Camel; C) The Battle of the Trench; D) The Battle of Badr
\\ \textit{Note:} The Battle of the Trench, which took place in 627 CE, was a pivotal moment for the early Muslim community, as it demonstrated their resilience against a coalition of tribes, ultimately leading to the weakening of their opponents and facilitating the subsequent peaceful conquest and conversion of Mecca in 630 CE.
\item \textbf{Answer:} C
\\ \textit{Options:} A) Shu and Set; B) Horus and Isis; C) Amun and Mut; D) Khonsu and Nekhbet
\\ \textit{Note:} Akhenaten sought to eliminate the worship of Amun and Mut in favor of the exclusive worship of the sun god Aten, reflecting his radical religious reforms that aimed to centralize worship and diminish the power of traditional priesthoods.
\item \textbf{Answer:} A
\\ \textit{Options:} A) Sita; B) Rama; C) Kali; D) Durga
\\ \textit{Note:} Sita is considered the ideal wife in Hindu tradition due to her unwavering devotion, loyalty, and virtue, exemplified in the epic Ramayana, where she endures hardships and remains faithful to her husband, Lord Rama.
\end{enumerate}

\section*{True / False}
\begin{enumerate}[label=\arabic*.]
\setcounter{enumi}{5}
\item \textbf{Answer:} True
\\ \textit{Options:} A) True; B) False
\\ \textit{Note:} The statement is true because the Conservative branch of Judaism, initiated by Zacharias Frankel, emphasizes a balance between tradition and modernity, which is encapsulated in the concept of 'Positive-Historical Judaism' that seeks to adapt Jewish practices while maintaining historical continuity.
\item \textbf{Answer:} True
\\ \textit{Options:} A) True; B) False
\\ \textit{Note:} Ganesha is indeed depicted with an elephant head in Hinduism and is revered as a symbol of wisdom and the remover of obstacles, making him a beloved figure associated with new beginnings.
\item \textbf{Answer:} False
\\ \textit{Options:} A) True; B) False
\\ \textit{Note:} The first Jaina temples actually appeared much earlier, around the 3 rd century BCE, during the time of the Maurya Empire, rather than in the second century CE.
\item \textbf{Answer:} True
\\ \textit{Options:} A) True; B) False
\\ \textit{Note:} The Gemarah is the compilation of rabbinical commentaries and discussions that expand upon the teachings of the Mishnah, making the statement true.
\item \textbf{Answer:} False
\\ \textit{Options:} A) True; B) False
\\ \textit{Note:} The central prayer in Jainism is actually known as the Navkar Mantra, not the Liberation Mantra.
\end{enumerate}

\section*{Long Form Response}
\begin{enumerate}[label=\arabic*.]
\setcounter{enumi}{10}
\item \textbf{Answer:} The ambrosial hours, which refer to the time between 3 a.m. and 6 a.m., hold great significance for Sikhs as they are considered the most spiritually potent time of the day. During these hours, Sikhs engage in meditation, prayer, and recitation of hymns, which deepens their connection to God and enhances their spiritual awareness. This practice is rooted in the belief that the early morning is a peaceful time, free from distractions, allowing for a focused and sincere devotion. By incorporating the ambrosial hours into their daily routines, Sikhs cultivate discipline and a sense of community, as many gather in gurdwaras for collective worship during this time. Thus, the ambrosial hours influence both personal spirituality and communal practices within Sikhism.
\\ \textit{Note:} correct because it accurately describes the ambrosial hours as a spiritually significant time for Sikhs, emphasizing their practices of meditation, prayer, and hymn recitation that foster a deeper connection to God in a distraction-free environment.
\item \textbf{Answer:} The Egyptian goddess Hathor was significant in ancient religion as she embodied love, beauty, motherhood, and joy, reflecting the values of Egyptian society. Often depicted as a cow or with cow horns, her imagery symbolizes fertility and nourishment, which were highly valued in a culture that relied on agriculture and livestock. Hathor's role as a goddess of music and dance also highlights the importance of celebration and community in Egyptian life. Additionally, her associations with the afterlife indicate the Egyptians' beliefs in life beyond death and the nurturing aspects of the divine. Overall, Hathor's depictions and attributes reveal a society that cherished familial bonds, prosperity, and spiritual fulfillment.
\\ \textit{Note:} correct because it effectively highlights Hathor's embodiment of essential values such as love, beauty, and motherhood, while illustrating her depictions as symbols of fertility and nourishment, which reflect the agricultural and communal priorities of ancient Egyptian society.
\end{enumerate}

\vfill
\noindent\textit{End of Answer Key}
\end{document}